\begin{textsample}{POS Dim 3 – pa – Score 8.00 – pa/pa\_inf\_5.txt}  \label{ex:f3_pos_280}
3.1.2.2 A Educação \textbf{Infantil} : importante etapa do processo de aprendizagem da \textbf{criança} Constituindo -se como a primeira etapa da Educação Básica, a Educação \textbf{Infantil} deve ser garantida em \textbf{creches}, para \textbf{crianças} de 0 a 3 anos, e em \textbf{pré-escolas}, para \textbf{crianças} de 4 e 5 anos, sob a responsabilidade prioritariamente dos poderes públicos municipais, integrando o Sistema Municipal de Ensino junto com Ensino Fundamental ( BRASIL, 1996)[12 ].
Dessa forma, com base na concepção de infância assumida neste documento curricular, reafirmamos o direito da \textbf{criança} ao atendimento educacional em consonância com a LDB n{\EmojiFont °} 9.394/1996.
A lei determina que a Educação \textbf{Infantil} tem como finalidade o desenvolvimento integral da \textbf{criança} até cinco anos de idade, em seus aspectos físicos, psicológicos, intelectuais e sociais, complementando a ação da família e da comunidade ( BRANDÃO, 2008 ).
No que se refere às duas formas de atendimento da Educação \textbf{Infantil}, é preciso analisar separadamente as faixas etárias de 0 a 3 anos e de 4 a 5 anos porque foram grupos tratados diferentemente, quer nos objetivos, quer nas demandas, quer nas instituições que atuam com essa etapa, sejam públicas ou privadas.
As preocupações com o atendimento de \textbf{crianças} da Educação \textbf{Infantil} devem se pautar na qualificação dos profissionais que atuam nessa etapa, além de se preocupar com o desenvolvimento dos programas e currículos, com a disponibilidade de \textbf{mobiliário}, equipamentos \textbf{lúdicos} e outros materiais pedagógicos adequados e necessários para os espaços.
É imperiosa a garantia de escolas de Educação \textbf{Infantil} às populações do campo, dos povos da floresta e dos rios, indígenas, quilombolas respeitando e garantindo assim essa etapa de ensino nos diferentes contextos amazônicos, privilegiando a constituição diversificada das \textbf{crianças} que neles habitam, respeitando, portanto, suas identidades, os seus aspectos socioculturais, étnico-raciais, de gênero, corporal, entre outros.
Apesar de a Educação \textbf{Infantil} e de o Ensino Fundamental serem etapas de escolarização diferentes, do ponto de vista da \textbf{criança} e da sua experiência não há fragmentação.
Nesse sentido, os professores e as instituições são os que muitas vezes se opõem e/ou fazem distinção desnecessária entre esses níveis de ensino, desconsiderando a \textbf{criança} e, consequentemente, negligenciado sua construção sócio-histórica e, sobretudo, sua experiência como sujeito cultural.
Questões relacionadas à alfabetização ou não na Educação \textbf{Infantil} ou no Ensino Fundamental ou como integrar esses dois níveis da Educação Básica, continuam recorrentes.
O importante é perceber que as \textbf{crianças} permanecem \textbf{crianças}, sejam na Educação \textbf{Infantil} ou no Ensino Fundamental e ainda, que esses níveis sejam indissociáveis, ou seja, que as \textbf{crianças} sejam oportunizadas de conhecimentos e \textbf{afetos}, saberes e valores, \textbf{cuidados} e atenção, seriedade e riso e, acima de tudo, ludicidade. [ 12 ] : A Constituição Federal de 1988, no capítulo VIII, Art. 227, estabelece o direito dos trabalhadores, pais e responsáveis, à educação de seus filhos e dependentes de 0 a 6 anos, além de considerar direito da própria \textbf{criança}.

% matched lemmas: afeto, creche, criança, cuidado, infantil, lúdico, mobiliário, pré-escola
\end{textsample}
